% ==============================================================
% CAPÍTULO 8 — CONCLUSÃO
% ==============================================================

\chapter{Conclusão}
\label{cap:conclusao}

\section{Síntese dos Achados}

Este relatório caracterizou o modelo de negócio, o ambiente competitivo e a estrutura organizacional da Capim Software, e analisou em profundidade seu processo decisório de inovação à luz de quatro frameworks da disciplina PRO3584. A tipologia de \textcite{salerno2015} mostrou que a empresa opera, na prática, mais de um tipo de processo de inovação simultaneamente -- tradicional para iniciativas incrementais, e de pausa ativa para iniciativas que dependem de maturação tecnológica ou de mercado --, mas sem reconhecer essa diferença na governança, aplicando um único critério trimestral de avaliação a projetos de natureza distinta. A cadeia de valor da inovação de \textcite{hansenbirkinshaw2007} indicou que o elo mais fraco do sistema da Capim é a geração de ideias: a empresa investiu em mecanismos sofisticados de conversão (seleção, métricas, financiamento por squad) e de difusão (progressão controlada por exposição a clínicas), mas não em um mecanismo equivalente de descoberta de oportunidades. A matriz de risco e retorno de \textcite{goffinmitchell2010} revelou um portfólio polarizado entre projetos incrementais de baixo risco e apostas radicais de alto risco, sem o estágio intermediário de oportunidades já validadas e de baixo risco que tipicamente preencheria o quadrante de \textit{Pearls}. Por fim, o modelo de ambidestria Descoberta--Incubação--Aceleração de \textcite{salernogomes2018} mostrou que a Capim já pratica uma forma de ambidestria estrutural ao isolar a Garagem~2 do restante da operação, mas não diferencia os critérios de decisão entre inovação incremental e radical -- lacuna que, segundo a literatura, é uma causa frequente de descontinuação precoce de iniciativas tecnologicamente promissoras.

Em resposta direta à devolutiva recebida sobre a versão anterior deste trabalho, conclui-se que a Capim tem, sim, um problema de geração de ideias -- não pela ausência de boas iniciativas, mas pela ausência de um mecanismo sistemático que as faça emergir, independentemente da disponibilidade pontual da liderança. A proposta de aprimoramento apresentada no Capítulo~\ref{cap:diagnostico-proposta} -- um hub de Descoberta, um \textit{learning plan} formal para iniciativas em incubação, a diferenciação explícita de critérios por tipo de processo, mecanismos formais de \textit{resourcing} e a revisão periódica do portfólio -- busca corrigir essa lacuna sem descartar os elementos do sistema atual que já se mostram alinhados às boas práticas discutidas em aula.

\section{Limitações}

Este estudo apresenta limitações que devem ser consideradas na interpretação de seus achados. Primeiro, trata-se de um estudo de caso único, baseado em uma única entrevista estruturada com um informante-chave (o PM de experimentos), o que restringe a possibilidade de generalização e introduz o risco de um relato unilateral do processo decisório -- outras lideranças ou squads poderiam descrever mecanismos de geração de ideias não mencionados nessa conversa específica. Segundo, a classificação das iniciativas da Capim segundo a tipologia de \textcite{salerno2015} e a matriz de \textcite{goffinmitchell2010} foi realizada pelos autores a partir do relato da entrevista, e não validada diretamente com a empresa -- um exercício de triangulação com a própria liderança da Capim fortaleceria a robustez do diagnóstico. Terceiro, o relatório não teve acesso a dados financeiros internos da empresa, o que limitou a análise de portfólio a uma leitura qualitativa, sem estimativas quantitativas de risco e retorno esperado por iniciativa.

\section{Trabalhos Futuros}

Como desdobramentos deste trabalho, sugere-se: (i) validar a tipologia de processos e a classificação de portfólio propostas neste relatório diretamente com a Diretoria de Produto e com os demais squads da Capim, ampliando a base de entrevistas; (ii) mensurar, ainda que qualitativamente, o número e a origem de ideias que efetivamente chegam à Garagem~2 ao longo de um período determinado, de modo a confirmar empiricamente o diagnóstico de fragilidade no elo de geração de ideias; e (iii) acompanhar, em um horizonte de um a dois trimestres, a evolução do projeto da CamilA caso a empresa venha a reavaliá-lo sob a lente de maturidade tecnológica proposta neste relatório, em vez do critério puramente financeiro/de adoção utilizado até o momento.
